相同味夸克散射比不同味过程多出一个不可回避的新结构：末态两个夸克完全相同，因此把两条外费米子线互换后得到的 $u$ 道图与原来的 $t$ 道图必须相干相加，而且 Fermi 统计给出相对负号。于是
\begin{equation}
\boxed{
\mathcal M(qq\to qq)=\mathcal M_t-\mathcal M_u,}
\label{eq:qcd-identical-qq-amplitude-dp95}
\end{equation}
其中 $\mathcal M_t$ 就是\eqnrefs{eq:qqprime-amplitude-dp10} 的同味版本，$\mathcal M_u$ 由 $p_3\leftrightarrow p_4$ 即 $t\leftrightarrow u$ 得到。对一般 $SU(N_c)$、无质量夸克，初态自旋和颜色平均后的结果为
\begin{equation}
\boxed{
\overline{|\mathcal M(qq\to qq)|^2}
=g_s^4\left\{
\frac{N_c^2-1}{2N_c^2}
\left[\frac{s^2+u^2}{t^2}+\frac{s^2+t^2}{u^2}\right]
-\frac{N_c^2-1}{N_c^3}\frac{s^2}{tu}
\right\}.}
\label{eq:qcd-identical-qq-m2-dp95}
\end{equation}
前两个平方项只是不同味结果及其交换像；第三项才是相同费米子交换的干涉。它同时依赖 Dirac 线的交换和颜色矩阵的非平凡收缩，不能通过把两个 Rutherford 型概率直接相加得到。物理 QCD 取 $N_c=3$ 后
\begin{equation*}
\overline{|\mathcal M|^2}
=g_s^4\left[
\frac49\frac{s^2+u^2}{t^2}
+\frac49\frac{s^2+t^2}{u^2}
-\frac{8}{27}\frac{s^2}{tu}
\right].
\end{equation*}

同一个四费米子解析振幅还把 $qq\to qq$ 与 $q\bar q\to q\bar q$ 连接起来。把一条外夸克从末态交叉到初态，相应四动量反号并把 $u$ 型自旋子换成 $v$ 型自旋子，Mandelstam 不变量随之置换。对相同味无质量 $q\bar q\to q\bar q$，树级同时存在 $s$ 道湮灭与 $t$ 道散射，结果为
\begin{equation}
\boxed{
\overline{|\mathcal M(q\bar q\to q\bar q)|^2}
=g_s^4\left\{
\frac{N_c^2-1}{2N_c^2}
\left[\frac{s^2+u^2}{t^2}+\frac{t^2+u^2}{s^2}\right]
-\frac{N_c^2-1}{N_c^3}\frac{u^2}{st}
\right\}.}
\label{eq:qcd-identical-qqbar-m2-dp95}
\end{equation}
这说明“交叉对称”不是把截面公式中的字母机械互换，而是先对解析振幅做外腿交叉，再重新进行物理通道的自旋、颜色平均；相同粒子统计和颜色干涉必须一起保留。

把\eqnrefs{eq:qcd-identical-qq-m2-dp95} 与 DP94 中的 $qg\to qg$、$gg\to gg$ 合在一起，QCD 的树级硬散射已经显出一个统一结构：运动学奇点由传播子通道决定，颜色系数由表示和生成元收缩决定，而完全相同外态还要求加入量子统计的交换干涉。后续的 PDF 因子化和喷注可观测量，就是把这些短程核嵌入强子初态并与软--共线长程结构分离。
